\cchapter{Data Environment}{data_environment}
\label{chap:data_environment}
An OpenMP \emph{data environment} is defined by a set of variables or objects
and their \emph{data-environment attributes}. Data-environment attributes can
be divided into \emph{data-sharing attributes} and \emph{data-mapping
attributes}. 

Many constructs (such as \kcode{parallel}, \kcode{simd}, \kcode{task}) 
accept clauses to control data-sharing attributes of referenced variables in
the construct, where data-sharing applies to whether the attribute of the
variable is \emph{shared} or \emph{private}, in addition to other special
operational characteristics of private (as indicated by the
\kcode{firstprivate}, \kcode{lastprivate}, \kcode{linear}, or \kcode{reduction}
clause).

Variables and objects in the data environment for a target device
(distinguished as a device data environment) have \emph{data-mapping
attributes} that are controlled by data-mapping constructs (such as
\kcode{target} or \kcode{target_data}), which determine the relationship of the
data on the host (the \emph{original} data) and the data on the device (the
\emph{corresponding} data).

\bigskip
DATA-SHARING ATTRIBUTES

Data-sharing attributes of variables can be classified as being \emph{predetermined},
\emph{explicitly determined} or \emph{implicitly determined}.

Certain variables and objects have predetermined attributes.  
A commonly found case is the loop iteration variable in associated loops 
of a \kcode{for} or \kcode{do} construct. It has a private data-sharing
attribute. Certain declarative directives can also be used to define variables
as having special predetermined data-sharing attributes including \emph{threadprivate},
\emph{groupprivate}, and \emph{device-local}. Variables with predetermined
data-sharing attributes cannot usually be listed in a data-sharing clause, but there
are some exceptions (mainly concerning loop iteration variables).

Variables with explicitly determined data-sharing attributes are those that are
referenced in a given construct and are listed in a data-sharing
clause on the construct. Some of the common data-sharing clauses are:
\kcode{shared}, \kcode{private}, \kcode{firstprivate}, \kcode{lastprivate}, 
\kcode{linear}, and \kcode{reduction}. % Are these all of them?

Variables with implicitly determined data-sharing attributes are those
that are referenced in a given construct, do not have predetermined
data-sharing attributes, and are not listed in a data-sharing
attribute clause of an enclosing construct.
For a complete list of variables and objects with predetermined and
implicitly determined attributes, please refer to the
\docref{Data-sharing Attribute Rules for Variables Referenced in a Construct}
subsection of the OpenMP Specification document.  

\bigskip
DATA-MAPPING ATTRIBUTES

A data-mapping attribute determines the manner in which a variable or object is
mapped from a data environment of a task (typically on the host device) to a
device data environment on a different device. The specification of list items
in a \kcode{map} clause is the main mechanism for controlling the data-mapping
attributes of data in a device data environment. These list items may include
variables, including array and structure elements, array sections, as well as
more general lvalue expressions in C/C++ (such as a dereferenced expression of
pointer type).

If a \kcode{map} clause is not explicitly specified for a variable that is
referenced in a \kcode{target} construct, that variable may still have an
\emph{implicit} data-mapping attribute (as if it had appeared in a \kcode{map}
clause). For example, the use of a declare target directive or
\kcode{defaultmap} clause can result in a variable having an implicit
data-mapping attribute. Additionally, list items that appear in certain
data-sharing clauses (e.g., \kcode{reduction}) on a compound target construct
can imply a data-mapping attribute. Also, non-scalar variables referenced
inside a \kcode{target} construct that do not otherwise have a predetermined or
explicit data-sharing or data-mapping attribute will typically be implicitly
mapped by default, in contrast to scalar variables which are typically given
an implicit firstprivate attribute (these default implicit attributes can be
changed with the use of the \kcode{defaultmap} clause). For a complete set of
rules for implicit data-mapping attributes, refer to the
\docref{Implicit Data-Mapping Attribute Rules}
subsection of the OpenMP Specification document.

The \kcode{map} clause can appear on data-mapping constructs (specifically,
\kcode{target}, \kcode{target_data}, \kcode{target_enter_data} and
\kcode{target_exit_data}).  The operations of creation and removal of corresponding
storage as well as assignment of the original list item values to the
corresponding list items may be complicated when the list item appears on
multiple constructs that are executed concurrently. To accomodate this, a
reference count is maintained to determine which of those operations are
needed. This can help ensure that corresponding storage is not removed on
completion of one construct while another construct that has mapped the same
data still requires it, as well as elide data transfers between devices in
cases where corresponding storage is not being created or removed (though this
can be overridden with use of modifiers such as \kcode{delete} or
\kcode{always}). Details of the \kcode{map} clause and reference count
operations are specified in the \docref{\kcode{map} Clause} subsection of the
OpenMP Specification document.


%===== Examples Sections =====
\input{data_environment/threadprivate}
\section{\kcode{groupprivate} Directive}
\label{sec:groupprivate}
\index{directives!groupprivate@\kcode{groupprivate}}
\index{groupprivate directive@\kcode{groupprivate} directive}

The \kcode{groupprivate} directive introduced in Specification 6.0
allows specified list items to be replicated such that each contention 
group will have its own uninitialized copy. The list item is shared among
threads of the contention group and does not exist 
outside the scope of the contention group. 

In the following example, the variable \ucode{x} is defined as a static 
variable and specified with the \kcode{groupprivate} data attribute in 
the function \ucode{foo}. Four teams created by the \kcode{teams} construct
execute the \kcode{parallel} region that calls the 
\ucode{foo} function. For each team the groupprivate variable \ucode{x} 
is created and is accessible for the group of tasks of the \kcode{parallel} region.

\cppexample[6.0]{groupprivate}{1}

\ffreeexample[6.0]{groupprivate}{1}

\input{data_environment/default_none}
\input{data_environment/private}
\input{data_environment/fort_loopvar}
\input{data_environment/fort_sp_common}
\input{data_environment/fort_sa_private}
\input{data_environment/fort_shared_var}
\input{data_environment/carrays_fpriv}
\input{data_environment/lastprivate}
\input{data_environment/reduction}
\input{data_environment/udr}
\input{data_environment/induction}
\input{data_environment/scan}
\input{data_environment/copyin}
\input{data_environment/copyprivate}
\input{data_environment/cpp_reference}
\input{data_environment/associate}

